\section{A complex neighborhood for the harmonic angular equation}
\label{ang:sec:kernel}

For real $a,b>0$ define, initially in the unit disk,
\begin{equation}
 F_{a,b}(z)=\sum_{n\ge2}\frac{H_{n-1}^{(b)}}{n^a}z^n,
 \qquad H_m^{(b)}=\sum_{j=1}^m j^{-b}.
 \label{ang:eq:F}
\end{equation}
This is the depth-two polylogarithm $\Li_{a,b}(z,1)$ when the
multiple-polylogarithm convention uses $n>m\ge1$ and the factor $z^n$.
The elementary inequality $1\le H_m^{(b)}\le m$ will be used
uniformly over all $b>0$.

\begin{lemma}[Positive kernel and exact moments]
\label{ang:lem:kernel}
On $\C\setminus[1,\infty)$,
\begin{equation}
 F_{a,b}(z)=z^2\int_{0<y<x<1}
 \frac{W_{a,b}(x,y)}{(1-xz)(1-yz)}\,dy\,dx,
 \quad
 W_{a,b}(x,y)=
 \frac{(-\log x)^{a-1}(\log(x/y))^{b-1}}{\Gamma(a)\Gamma(b)}.
 \label{ang:eq:kernel}
\end{equation}
For nonnegative integers $j,k$ its moments are
\begin{equation}
 \int_{0<y<x<1}W_{a,b}(x,y)x^jy^k\,dy\,dx
   =\frac1{(j+k+2)^a(k+1)^b}.
 \label{ang:eq:moments}
\end{equation}
\end{lemma}
\begin{proof}
Set $x=e^{-t}$ and $y=e^{-t-v}$. The Jacobian is $xy$, and the
integral in \eqref{ang:eq:moments} separates into Gamma integrals
with decay rates $j+k+2$ and $k+1$. Expanding both geometric
denominators in \eqref{ang:eq:kernel}, the coefficient of $z^n$
is
\[
 \sum_{k=0}^{n-2}\frac1{n^a(k+1)^b}
 =\frac{H_{n-1}^{(b)}}{n^a}.
\]
Absolute convergence justifies this computation for $|z|<1$.
On every compact subset of the slit plane the denominators are bounded
away from zero, and $W$ has finite mass $2^{-a}$. Dominated
convergence and differentiation on compact sets give the holomorphic
continuation. This is the kernel of \cite{RepoThreshold}, with its
moments made explicit for later estimates.
\end{proof}

For $a\ge6$, $b>0$, and $0<\rho\le1$, let $\theta_{a,b}(\rho)$ be the unique zero
in $(0,\pi)$ of $\operatorname{Im}F_{a,b}(\rho e^{i\theta})$.
In the range needed here its existence and simplicity follow from
the global angular theorem of \cite{RepoThreshold}; alternatively,
the construction below independently gives the root near $\pi/2$.
Write
\[
 \delta_{a,b}(\rho)=\frac\pi2-\theta_{a,b}(\rho).
\]
The Gaussian sign in the same source puts this displacement in
$(0,\pi/2)$. We shall also see its local positivity directly.

Set
\begin{align}
 R_{a,b}(z)&=F_{a,b}(z)-2^{-a}z^2,\notag\\
 z_+(v)&=i\rho e^{-iv},\qquad z_-(v)=-i\rho e^{iv},\notag\\
 S(v)&=\frac{2^a}{2i\rho^2}
       \bigl(R_{a,b}(z_+(v))-R_{a,b}(z_-(v))\bigr),
 \label{ang:eq:S}\\
 B(v)&=\frac{v}{\sin2v},\qquad B(0)=\frac12,
 \qquad f(v)=-B(v)S(v).
 \label{ang:eq:f}
\end{align}
These definitions are holomorphic; the difference in \eqref{ang:eq:S}
is used instead of taking an imaginary part at complex $v$.
For real $v$ it equals the normalized imaginary tail. The angular
equation is exactly
\begin{equation}
 \sin2v+S(v)=0.
 \label{ang:eq:angular}
\end{equation}

\begin{lemma}[Uniform analytic tail bound]
\label{ang:lem:norm}
Let $0<r<\pi/2$ and $c=\cos r$. On the disk $|v|\le r$,
\begin{equation}
 |S(v)|\le
 \rho e^{3r}\bigl(c^{-2}+2^{-b}c^{-1}\bigr)(2/3)^a.
 \label{ang:eq:norm}
\end{equation}
At $r=1/4$ this implies
\begin{equation}
 \sup_{|v|\le1/4}|f(v)|<\frac14\eta_{a,\rho},
 \qquad \eta_{a,\rho}=10\rho(2/3)^a.
 \label{ang:eq:eta}
\end{equation}
In particular $a\ge6$ gives $\eta_{a,\rho}\le640/729<1$.
\end{lemma}
\begin{proof}
If $z$ lies in either sector of half-width $r$ about $i$ or $-i$,
then $|1-tz|\ge\cos r=c$ for every $t\ge0$: this is the
distance from $1$ to a ray whose angle is at least $\pi/2-r$ in
absolute value. The elementary identity
\[
 \frac1{(1-A)(1-C)}-1
 =\frac{A}{(1-A)(1-C)}+\frac{C}{1-C}
\]
and the moments of $Wx$ and $Wy$ give
\[
 |R_{a,b}(z)|
 \le |z|^3 3^{-a}\bigl(c^{-2}+2^{-b}c^{-1}\bigr).
\]
For $|v|\le r$, both $z_\pm(v)$ lie in these sectors and have
modulus at most $\rho e^r$. Substitution proves
\eqref{ang:eq:norm}.

For completeness the constants can be bounded rationally:
\[
 \cos(1/4)>31/32,\qquad e^{3/4}<17/8,
 \qquad \sup_{|v|\le1/4}|B(v)|<8/15.
\]
The last inequality follows from
$|\sin(2v)/(2v)-1|\le\sinh(1/2)/(1/2)-1<1/16$.
For that strict bound the positive tail is at most
$(1/24)/(1-1/80)=10/237<1/16$.
Taylor's formula and a geometric bound on its positive tail give
$e^{3/4}<18427/8704<17/8$; the cosine bound follows from
$\cos x>1-x^2/2$. Consequently
\[
 4\sup|B|e^{3/4}(c^{-2}+c^{-1})
 <\frac{45696}{4805}<10.
\]
This proves \eqref{ang:eq:eta}, including its strictness.
\end{proof}

\section{A convergent summation of the angular multiset transseries}
\label{ang:sec:summation}

The scalar parameter used below scales the complete perturbation
before taking boundary values. It is not an independent variation of
infinitely many Fourier coefficients. Define the coupled equation
\begin{equation}
 \delta(\tau)=\tau f(\delta(\tau)),\qquad \delta(0)=0.
 \label{ang:eq:coupling}
\end{equation}

\begin{theorem}[Convergent angular reconstruction]
\label{ang:thm:sum}
For real $a\ge6$, $b>0$, and $0<\rho\le1$, put
\begin{equation}
 D_k(a,b;\rho)=\frac1k[v^{k-1}]f(v)^k
 =\frac{(-1)^k}{k}[v^{k-1}]B(v)^kS(v)^k.
 \label{ang:eq:D}
\end{equation}
Then
\begin{equation}
 \delta_{a,b}(\rho)=\sum_{k\ge1}D_k(a,b;\rho)
 \label{ang:eq:sum}
\end{equation}
converges absolutely. More precisely,
\begin{align}
 |D_k(a,b;\rho)|&\le\frac{\eta_{a,\rho}^k}{4k},
 \label{ang:eq:D-bound}\\
 \left|\delta_{a,b}(\rho)-\sum_{k=1}^KD_k(a,b;\rho)\right|
 &\le\frac{\eta_{a,\rho}^{K+1}}
 {4(K+1)(1-\eta_{a,\rho})},\qquad K\ge0.
 \label{ang:eq:sum-error}
\end{align}
The solution of \eqref{ang:eq:coupling} is holomorphic on
$|\tau|<\eta_{a,\rho}^{-1}$, which contains $\tau=1$.
\end{theorem}
\begin{proof}
On $|v|=1/4$, Lemma~\ref{ang:lem:norm} gives
$|\tau f(v)|<|v|$ when $|\tau|\eta_{a,\rho}<1$.
Rouch\'e's theorem shows that $v-\tau f(v)$ has exactly one
zero, counted with multiplicity, in the disk. It is therefore simple.
The implicit function theorem and uniqueness patch its local branches
into one holomorphic function of $\tau$ on the indicated disk.
Lagrange inversion gives \eqref{ang:eq:D}.
Applying Cauchy's coefficient estimate to $f(v)^k$ on the fixed
circle gives
\[
 |D_k|\le\frac1k(1/4)^{1-k}
              (\eta_{a,\rho}/4)^k
 =\frac{\eta_{a,\rho}^k}{4k}.
\]

At real $\tau$ the zero is real by conjugation symmetry. Moreover,
$S(0)<0$. Indeed the imaginary part of the positive kernel at
$z=i\rho$ is strictly negative: $z^2=-\rho^2$, while
the imaginary part of $((1-i\rho x)(1-i\rho y))^{-1}$ is
strictly positive for $0<y<x<1$. The removed quadratic term is
real. Thus $f(0)>0$, the zero starts positive for small positive
$\tau$, and it cannot cross zero for any $\tau>0$.
At $\tau=1$ it is in $(0,1/4)$ and solves the actual angular
equation. Global angular uniqueness identifies it with the stated
$\delta_{a,b}(\rho)$. Finally summing
$\sum_{k>K}\eta^k/k\le\eta^{K+1}/((K+1)(1-\eta))$
proves the error bound.
\end{proof}

\subsection{The exact summation order}

A multiset of Fourier modes is a finitely supported sequence
$\boldsymbol m=(m_3,m_4,\ldots)$ of nonnegative integers. Set
\[
 k=|\boldsymbol m|=\sum m_n,\quad
 \Lambda(\boldsymbol m)=\prod_{n\ge3}(n/2)^{m_n},\quad
 R(\boldsymbol m)=\sum_{n\ge3}(n-2)m_n,
 \quad h_n=H_{n-1}^{(b)}.
\]
The source coefficient is
\begin{equation}
 C(\boldsymbol m)=\frac{(-1)^k(k-1)!}{\prod m_n!}
 [v^{k-1}]B(v)^k
 \prod_{n\ge3}\sin(n\pi/2-nv)^{m_n}.
 \label{ang:eq:multiset}
\end{equation}
The source proves that the sum with $\Lambda<B_0$, for fixed
$B_0>1$, approximates the displacement with error
$O_{B_0}(\rho B_0^{-a})$ as $a\to\infty$~\cite{RepoAngular}.
The support is locally finite in $\Lambda$: every participating
$n$ is less than $2B_0$ and $k<\log B_0/\log(3/2)$.

\begin{theorem}[Degree blocks and their Abel boundary values]
\label{ang:thm:Abel}
For $0<\rho<1$ and fixed $k$, the following inner sum converges
absolutely:
\begin{equation}
 D_k(a,b;\rho)=\sum_{|\boldsymbol m|=k}
 C(\boldsymbol m)\rho^{R(\boldsymbol m)}
 \prod h_n^{m_n}\,\Lambda(\boldsymbol m)^{-a}.
 \label{ang:eq:blocks}
\end{equation}
For $a\ge6$, $b>0$, its radial boundary values satisfy
\begin{align}
 D_k(a,b;1)&=\lim_{\rho\uparrow1}
  \sum_{|\boldsymbol m|=k}
  C(\boldsymbol m)\rho^{R(\boldsymbol m)}
  \prod h_n^{m_n}\,\Lambda(\boldsymbol m)^{-a},\notag\\
 \delta_{a,b}(1)&=\sum_{k\ge1}D_k(a,b;1).
 \label{ang:eq:Abel}
\end{align}
The final sum is ordinary and absolutely convergent, with
\eqref{ang:eq:sum-error} at $\rho=1$.
\end{theorem}
\begin{proof}
For subunit $\rho$, on a sufficiently small complex disk
$|v|<\epsilon<-\log\rho$, the Fourier expansion
\[
 S(v)=\sum_{n\ge3}h_n\rho^{n-2}(2/n)^a
                    \sin(n\pi/2-nv)
\]
and every fixed number of its derivatives converge absolutely.
For fixed $k$ one may expand its $k$th power and extract the
coefficient of degree $k-1$ absolutely. The multinomial coefficient
and the factor $1/k$ give exactly \eqref{ang:eq:multiset}.

For $\rho\uparrow1$, the kernel formula gives uniform convergence
of $S$ on $|v|\le1/4$: the denominators have a common nonzero
lower bound and the numerator is dominated by a constant times the
integrable weight. Each fixed Taylor coefficient in
\eqref{ang:eq:D} therefore converges. The summable majorant
$\eta_{a,1}^k/(4k)$ is independent of $\rho$, so dominated
convergence also interchanges the radial limit with the outer degree
sum. This proves all the assertions with the displayed order of
operations.
\end{proof}

\begin{remark}[Which data the summation uses]
This is a summation of the canonical multiset family, retaining both
$|\boldsymbol m|$ and $R(\boldsymbol m)$. Equal rational
$\Lambda$ can arise with different degrees and different radial
weights. The prescription is therefore not asserted to be a
summation functional depending only on the already grouped scalar
coefficients $\widehat C(\Lambda;b)$. All arithmetic collisions
are retained in \eqref{ang:eq:blocks}; they are not discarded.
\end{remark}

\subsection{Why divergence persists in the original ordering}

For clarity, the divergence proof in \cite{RepoFormal} is recalled
in a short form. It also extends verbatim from integer $a$ to every
fixed real $a>0$. Choose an integer $s$ with $2s>a$. The multiset
with $2s$ copies of mode $3$ and one copy of an odd prime $P\ne3$
has exponential base $q_P=2^{2s+1}/(3^{2s}P)$. This base has a
unique multiset representation. If another representation has $k$
factors, then
\[
 \prod_{j=1}^k n_j=2^{k-(2s+1)}3^{2s}P.
\]
Integrality first forces $k\ge2s+1$. Counting prime factors then
forces every $n_j$ to be prime. Since $n_j\ge3$, no factor two
is possible; hence $k=2s+1$ and the representation is unique.

The coefficient extraction contains
\[
 [v^{2s}]B(v)^{2s+1}\cos^{2s}(3v)\cos(Pv)
 =\frac{(-1)^sP^{2s}}{2^{2s+1}(2s)!}+O_s(P^{2s-2}).
\]
Its term after multiplication by $q_P^a$ has magnitude a positive
constant times
$H_{P-1}^{(b)}P^{2s-a}(1+O(P^{-2}))$, which tends to infinity.
Thus even the terms grouped by equal bases fail to tend to zero.
In \eqref{ang:eq:Abel}, by contrast, the inner boundary blocks
are Abel limits, not ordinary sums. Absolute convergence of the
outer degree series does not assert absolute convergence of the
flattened multiset family, even for every subunit $\rho$.

\section{Finite polylogarithm coordinates for every degree}
\label{ang:sec:periods}

Let $S_j=S^{(j)}(0)$. Differentiating the holomorphic definition
\eqref{ang:eq:S}, rather than a boundary Fourier series, yields
\begin{equation}
 S_j=\frac{2^a}{2i\rho^2}
 \left((-i)^jF_{a-j,b}(i\rho)-i^jF_{a-j,b}(-i\rho)\right)
 -2^j\sin(j\pi/2).
 \label{ang:eq:jets}
\end{equation}
Lower outer orders here are defined by the analytic identity
$z\partial_zF_{s,b}=F_{s-1,b}$. For any finite number of
derivatives one can start with a positive outer order and differentiate
its slit-plane continuation, so this notation is unambiguous.

The first four degree blocks are
\begin{align}
 D_1&=-\frac{S_0}{2},\qquad
 D_2=\frac{S_0S_1}{4},\notag\\
 D_3&=-\frac{S_0^3}{12}-\frac{S_0^2S_2}{16}
                         -\frac{S_0S_1^2}{8},\notag\\
 D_4&=\frac{S_0^3S_1}{6}+\frac{S_0^3S_3}{96}
       +\frac{3S_0^2S_1S_2}{32}+\frac{S_0S_1^3}{16}.
 \label{ang:eq:first-D}
\end{align}
More generally, if $B(v)=\sum\beta_jv^j$,
$S(v)=\sum s_jv^j$, and $f_j=-\sum_{\ell=0}^j\beta_\ell
s_{j-\ell}$, then
\begin{equation}
 D_k=\frac1k\sum_{j_1+\cdots+j_k=k-1}
                 f_{j_1}\cdots f_{j_k}.
 \label{ang:eq:finite-composition}
\end{equation}
This is a finite polynomial formula, evaluated by truncated polynomial
multiplication in the accompanying program.

\begin{proposition}[Reduction at nonpositive outer order]
\label{ang:prop:lowering}
For every integer $m\ge0$ and real $b>0$,
\begin{equation}
 F_{-m,b}(z)=\sum_{j=0}^m\binom mj
             \Li_{j-m}(z)\Li_{b-j}(z)
 \label{ang:eq:lowering}
\end{equation}
on the common slit-plane continuation.
\end{proposition}
\begin{proof}
Summing first over $n>j$ in the unit disk gives
$F_{0,b}(z)=\Li_0(z)\Li_b(z)$, where
$\Li_0(z)=z/(1-z)$. Apply $(z\partial_z)^m$, use
Leibniz' rule and $z\partial_z\Li_s=\Li_{s-1}$, and
then analytically continue. Each factor with nonpositive integral
order is rational. The identity itself is elementary; its role is
to close the coordinate algebra of the convergent reconstruction.
\end{proof}

\begin{corollary}[A fixed period algebra]
\label{ang:cor:algebra}
For integers $a\ge6$, $b\ge1$, every $D_k(a,b;1)$ belongs to
the same finitely generated real algebra
\begin{equation}
 \Q\left[
 \begin{array}{ll}
 \operatorname{Re}F_{r,b}(i),\operatorname{Im}F_{r,b}(i),&1\le r\le a,\\
 \operatorname{Re}\Li_j(i),\operatorname{Im}\Li_j(i),&1\le j\le b
 \end{array}\right].
 \label{ang:eq:algebra}
\end{equation}
Increasing the nonlinear degree introduces no additional depth-two
generators or classical-polylogarithm generators.
\end{corollary}
\begin{proof}
Use \eqref{ang:eq:jets}, \eqref{ang:eq:finite-composition}, and
\eqref{ang:eq:lowering}. All nonpositive integral-order
polylogarithms at $i$ belong to $\Q(i)$, and values at $-i$ are
complex conjugates. The operations required are finite sums and
products with rational coefficients. This proves membership, not
minimality or algebraic independence of the generating set.
\end{proof}

\section{A moving finite truncation with Gaussian error}
\label{ang:sec:moving}

The convergent degree reconstruction still uses analytically evaluated
polylogarithms in each block. It is useful to have a completely finite
alternative involving only a finite number of harmonic numbers and
trigonometric Taylor coefficients. Let $M\ge3$ be an integer and put
\begin{align}
 S_M(v)&=\sum_{n=3}^M H_{n-1}^{(b)}\rho^{n-2}(2/n)^a
                         \sin(n\pi/2-nv),\notag\\
 f_M(v)&=-B(v)S_M(v),\qquad
 D_{k,M}=\frac1k[v^{k-1}]f_M(v)^k,\notag\\
 P_{M,K}(a,b;\rho)&=\sum_{k=1}^K D_{k,M}.
 \label{ang:eq:finite-cutoff}
\end{align}
This keeps modes at most $M$ and total degree at most $K$. It is
a finite subfamily of the original multiset formula, with every
collision inside that subfamily combined automatically.

\begin{theorem}[Explicit finite approximation bound]
\label{ang:thm:finite}
Let $a\ge8$, $b>0$, $0<\rho\le1$, $M\ge3$, and $K\ge0$.
If
\begin{equation}
 \kappa=14M\rho(2/3)^a\le\frac12,
 \label{ang:eq:kappa}
\end{equation}
then
\begin{equation}
 \left|\delta_{a,b}(\rho)-P_{M,K}(a,b;\rho)\right|
 \le \frac{\rho2^aM^{2-a}}{a-2}
     +\frac{\kappa^{K+1}}{M(K+1)(1-\kappa)}.
 \label{ang:eq:finite-error}
\end{equation}
The two displayed terms control frequency removal and degree removal,
respectively.
\end{theorem}
\begin{proof}
Use the disk $|v|\le1/M$. Here $|B(v)|\le1$ and
$|\sin(n\pi/2-nv)|\le e^{n/M}\le e$ for $n\le M$.
For $a\ge4$,
\[
 \sum_{n\ge3}(n-1)(2/n)^a
 \le (2/3)^a\,81\left(\zeta(3)-\zeta(4)-\frac1{16}\right)
 <5(2/3)^a.
\]
Since $5e<14$, it follows that
$\sup|f_M|<14\rho(2/3)^a=\kappa/M$.
The same Rouch\'e and Lagrange argument on radius $1/M$ gives
a unique real small root $\delta_M$ and
\[
 |D_{k,M}|\le\frac{\kappa^k}{Mk},\qquad
 \left|\delta_M-\sum_{k=1}^K D_{k,M}\right|
 \le\frac{\kappa^{K+1}}{M(K+1)(1-\kappa)}.
\]

The actual root is also in this real interval. At its fixed point,
Lemma~\ref{ang:lem:norm} gives
$|\delta_{a,b}|\le(10/4)\rho(2/3)^a<1/M$ under
\eqref{ang:eq:kappa}. On $|v|\le1/M$ the real derivative of
the finite equation $E_M(v)=\sin2v+S_M(v)$ satisfies
\begin{align*}
 E_M'(v)
 &\ge2\cos(2/M)-\sum_{n\ge3}n(n-1)(2/n)^a\\
 &\ge2\cos(2/3)-14(2/3)^a>1.
\end{align*}
For the second line, comparison at exponent five uses
\[
 \sum_{n\ge3}n(n-1)(3/n)^5
 =243\left(\zeta(3)-\zeta(4)-\frac1{16}\right)<14.
\]
This inequality has a finite rational certificate: sum
$(n-1)/n^4$ from $3$ through $100$ and bound the decreasing
tail by $1/(2\cdot100^2)-1/(3\cdot100^3)$. Multiplying by
$243$ gives a rational number less than $14$. Finally,
$2\cos(2/3)\ge14/9$ and
$14/9-14(2/3)^8=6622/6561>1$.

At the actual root the residual of the finite equation is its removed
Fourier tail, whose magnitude is at most
\[
 \rho2^a\sum_{n>M}n^{1-a}
 \le\frac{\rho2^aM^{2-a}}{a-2}.
\]
The derivative lower bound and the mean value theorem control
$|\delta_{a,b}-\delta_M|$ by this residual. Adding the degree
error proves \eqref{ang:eq:finite-error}.
\end{proof}

\begin{theorem}[An attainable Gaussian error scale]
\label{ang:thm:gaussian}
Put $\ell=\log(3/2)$. For every fixed $0<\lambda<\ell$ choose
\begin{equation}
 M(a)=\lfloor e^{\lambda a}\rfloor,\qquad
 K(a)=\left\lceil\frac{\lambda a}{\ell-\lambda}\right\rceil.
 \label{ang:eq:moving-choice}
\end{equation}
There are constants $a_\lambda,C_\lambda>0$, independent of
$b>0$ and $0<\rho\le1$, such that for $a\ge a_\lambda$,
\begin{equation}
 |\delta_{a,b}(\rho)-P_{M(a),K(a)}(a,b;\rho)|
 \le \exp(-\lambda a^2+C_\lambda a).
 \label{ang:eq:gaussian}
\end{equation}
\end{theorem}
\begin{proof}
We have $\log M=\lambda a+O(e^{-\lambda a})$ and
$\kappa\le14e^{-(\ell-\lambda)a}$, so the hypotheses of
Theorem~\ref{ang:thm:finite} hold eventually, uniformly in
$b,\rho$. The logarithm of its frequency bound is at most
$-\lambda a^2+(2\lambda+\log2)a+O(\log a)$.
For the degree bound use
$K+1\ge\lambda a/(\ell-\lambda)$ and
$\log\kappa\le\log14-(\ell-\lambda)a$.
Its logarithm is likewise at most $-\lambda a^2+O_\lambda(a)$.
Increasing $C_\lambda$ absorbs their sum and all fixed thresholds.
\end{proof}

A simple choice is $\lambda=\ell/2$, giving
$M=\lfloor(3/2)^{a/2}\rfloor$ and $K=\lceil a\rceil$.
The theorem measures error, not computational optimality: $M$ grows
exponentially with $a$. It is a proved partial answer to the moving
truncation question, while the sharp error of a single
exponential-base cutoff remains open.

\subsection{An exact complex frequency remainder}

The positive kernel gives another useful finite identity. Put
$h_d(x,y)=\sum_{j=0}^d x^{d-j}y^j$. Direct multiplication gives
\begin{equation}
 \frac1{(1-xz)(1-yz)}-\sum_{d=0}^{M-3}h_d(x,y)z^d
 =\frac{z^{M-2}(h_{M-2}(x,y)-zxyh_{M-3}(x,y))}
 {(1-xz)(1-yz)}.
 \label{ang:eq:exact-tail}
\end{equation}
In the sectors of Lemma~\ref{ang:lem:norm}, integration and the
exact moments therefore imply
\begin{align}
 \left|F_{a,b}(z)-\sum_{n=2}^{M-1}\frac{H_{n-1}^{(b)}}{n^a}z^n\right|
 &\le\frac{|z|^M}{c^2}
 \left(\frac{H_{M-1}^{(b)}}{M^a}
 +|z|\frac{H_{M-1}^{(b)}-1}{(M+1)^a}\right).
 \label{ang:eq:complex-tail}
\end{align}
Unlike an estimate obtained by summing the defining Fourier tail,
this remains valid for $|z|>1$ in those sectors. At fixed complex
radius $r$, its normalized angular version is
$O_r(\rho^{M-2}M e^{Mr}(2/M)^a)$. Taking $M$ proportional
to $a/r$ yields $\exp(-a\log a+O_r(a))$. The stronger Gaussian
construction above instead shrinks the analytic disk with $M$;
both estimates describe their precise domains without assuming
convergence of the original infinite transseries.
